
\documentclass{article}



\usepackage{graphicx,comment,framed}
\usepackage{roundbox}
\usepackage{fancybox}
\usepackage{tikz}
\usepackage{color}
\usepackage[hidelinks]{hyperref}
\usepackage{framed}
\usepackage{amsthm,amssymb,amsmath}
%\usepackage[colorlinks,linkcolor=blue,citecolor=blue]{hyperref}
\definecolor{shadecolor}{cmyk}{0,0,0,0}
\usepackage{listings}
\usepackage{xepersian}
\usepackage[noend]{algpseudocode}


%--------------------- page settings ----------------------

\settextfont[Scale=1.1]{XB Niloofar}
\setdigitfont[Scale=1.1]{XB Niloofar}
\defpersianfont\sayeh[Scale=1.1]{XB Niloofar}
\addtolength{\textheight}{3.2cm}
\addtolength{\topmargin}{-22mm}
\addtolength{\textwidth}{3cm}
\addtolength{\oddsidemargin}{-1.5cm}


%------------------------ Environments ------------------------------------

\newtheorem{قضیه}{قضیه}
\newtheorem{لم}{لم}
\newtheorem{مشاهده}{مشاهده}
\newtheorem{تعریف}{تعریف}


%-------------------------- Notations ------------------------------------

\newcommand{\IR}{\ensuremath{\mathbb{R}}} 
\newcommand{\IZ}{\ensuremath{\mathbb{Z}}} 
\newcommand{\IN}{\ensuremath{\mathbb{N}}} 
\newcommand{\IS}{\ensuremath{\mathbb{S}}} 
\newcommand{\IC}{\ensuremath{\mathbb{C}}} 
\newcommand{\IB}{\ensuremath{\mathbb{B}}} 

\newcommand{\bR}{\mathbb{R}}
\newcommand{\cB}{\mathcal{B}}
\newcommand{\cO}{\mathcal{O}}
\newcommand{\cG}{\mathcal{G}}
\newcommand{\rM}{\mathrm{M}}
\newcommand{\rC}{\mathrm{C}}
\newcommand{\rV}{\mathrm{V}}

\newcommand{\lee}{\leqslant}
\newcommand{\gee}{\geqslant}
\newcommand{\ceil}[1]{{\left\lceil{#1}\right\rceil}}
\newcommand{\floor}[1]{{\left\lfloor{#1}\right\rfloor}}
\newcommand{\prob}[1]{{\mbox{\tt Pr}[#1]}}
\newcommand{\set}[1]{{\{ #1 \}}}
\newcommand{\seq}[1]{{\left< #1 \right>}}
\newcommand{\provided}{\,|\,}
\newcommand{\poly}{\mbox{\rm poly}}
\newcommand{\polylog}{\mbox{\rm \scriptsize polylog}\,}
\newcommand{\comb}[2] {\left(\!\!\begin{array}{c}{#1}\\{#2}\end{array}\!\!\right)}




\newcounter{probcnt}
\newcommand{\مسئله}[1]{\stepcounter{probcnt}\bigskip\bigskip{
 	\large \bf مسئله‌ی \arabic{probcnt}$\mbox{\bf{.}}$ \ #1} \bigskip}

\newcommand{\fqed}[1]{\leavevmode\unskip\nobreak\quad\hspace*{\fill}{\ensuremath{#1}}}

\newenvironment{اثبات}
	{\begin{trivlist}\item[\hskip\labelsep{\em اثبات.}]}
	{\fqed{\square}\end{trivlist}}

\newenvironment{حل}
	{\begin{trivlist}\item[\hskip\labelsep{\bf حل.}]}
	{\fqed{\blacktriangleright}\end{trivlist}}

\ifdefined\hidesols
	\newsavebox{\trashcan} % uncomment the following line to hide solutions
	\renewenvironment{حل}{\begin{latin}\begin{lrbox}{\trashcan}}{\end{lrbox}\end{latin}}
\fi


%------------------------- Header -----------------------------

%------------------------- Header -----------------------------

\newcommand{\سربرگ}[3]{
	\parindent=0em
	\begin{shaded}
		
		\rightline{ 
			\makebox[6em][c]{
				\includegraphics[height=1.6cm]{aut.png}
		}}
		\vspace{-.2em}
		{\scriptsize\bf دانشکده‌ی مهندسی کامپیوتر}
		\hfill {\small
			مدرس: دکتر امیرمزلقانی  \ 
		}\\[-5em]
		\leftline{\hfill\Large\bf 
			کاربرد های جبر خطی
		}\\[.7em]
		\leftline{\hfill\bf 
			نیم‌سال دوم ۹۶-۹۷
		}\\[1.7em]
		\hrule height .12em
		
		\normalsize
		\vspace{1mm} #1
		\hfill \small  #3
		\vspace{1mm} 
		\hrule height .1em
		
		\vspace{-0.5em} 
		\hfill {\sayeh\large #2} \hfill
	\end{shaded}
	\begin{large}

توجه:
\end{large}
\\
\\
\begin{itemize}
\item
این تمرین از مباحث مربوط به فصل5 ،6و7 طراحی شده است که شامل 15سوال اجباری و 4 سوال امتیازی است که نمره سوال های امتیازی فقط به نمرات تمرین شما کمک می کند.
\item 
در برخی سوالات چند قسمتی از شما خواسته شده است به تعدادی از قسمت ها به اختیار پاسخ دهید ،پاسخگویی به بیش از مقدار تعیین شده نمره اضافی ندارد ،از بقیه قسمت ها می توانید برای تمرین بیشتر استفاده کنید.
\item 
در پایان نیز تعدادی سوال برای تمرین بیشتر در نظر گرفته شده است ،که به حل آن ها نمره ای تعلق نمی گیرد صرفا برای تمرین توسط خود شما در نظر گرفته شده است. 
\item 
اگه سوالی داشتین از طریق
\begin{latin}
\href{mailto:aut.la2018@gmail.com}{\[aut.la2018@gmail.com\]}
\end{latin}
 حتما بپرسید.

\item پاسخ های تمرین را در قالب یک فایل به صورت الگوی زیر آپلود کنید.
\begin{latin}
9531000\_steve\_McManaman\_HW4.pdf
\end{latin}
\item \color{red}
مهلت تحویل جمعه  1تیر 1397 ساعت 23:54:59
\end{itemize}

}




%\settextfont{Yas}
\begin{document}
\سربرگ{تمرین اول}{}{}
\clearpage
\مسئله{}  ماتریس  های زیر متعلق ماتریس افزوده سه دستگاه معادله خطی است،در هرمرحله پس از مشخص  کردن جایگاه(درایه) و ستون محوری و با استفاده از روش حذفی گاوس جردن ماتریس ها را به شکل کاهش یافته سطری در بیاورید و سپس در مورد جواب دستگاه ها بحث کنید.(در صورت داشتن داشتن جواب عمومی،جواب ها را به صورت پارامتریک بنویسید.)
\\
الف)
${\begin{bmatrix}
		1&2&3&4&1\\
		0&1&2&3&1\\
		0&0&0&1&1\\
		0&1&0&1&1
		\end{bmatrix}}$
	\qquad
	ب)
$\begin{bmatrix}
1&-7&0&6&5\\
0&0&1&-2&-3\\
-1&7&-4&2&7
\end{bmatrix}
$
\qquad
ج)
$\begin{bmatrix}
3&-4&2&0\\
-9&12&-6&0\\
-6&8&-4&0
\end{bmatrix}$

\مسئله{}
مقدار 
		$\lambda$
		را طوری تعیین کنید که 
		\\
		الف) دستگاه معادلات زیر جواب ناصفر داشته باشد:
	\begin{equation*}
	 \left\{
	\begin{array}{rl}
	 \lambda x+3y=0\\
	 2x+4y=0
	\end{array} \right.
	\end{equation*}
	
	(بیش از آنکه جواب نهایی شما مهم باشد روش حل و نوع نگاه شما به آن مهم است.روشی که قابل تعمیم به مسائل دیگر باشد)
	\\
	ب) دستگاه معادلات زیر فقط جواب صفر داشته باشد :
		\begin{equation*}
	\left\{
	\begin{array}{rl}
	\lambda x_1+x_2+x_3=0\\
	x_1+\lambda x_2+x_3=0\\
	x_1+x_2+x_3=0
	\end{array} \right.
	\end{equation*}
	
\مسئله{} درستی یا نادرستی گزاره های زیر را مشخص کنید و درصورت نادرست بودن مثال نقض ارائه دهید و در صورت درست بودن آن را اثبات کنید:
\\
الف) اگر 
$\forall i\ \ v_i\in {\Bbb R}^n$
و 
$\{v_1,v_2,\ldots,v_n \}$
یک مجموعه وابسته خطی باشد ،هریک از $v_i$ ها را می توان به صورت یک ترکیب خطی از بقیه اعضا نوشت.
\\
ب) اگر $s_1 $ و $s_2 $  زیر مجموعه هایی از بردار های 
${\Bbb R}^n$
باشند که 
$span(s_1)=span(s_2)$
آنگاه 
$s_1=s_2$
.
\\
ج) اگر
 $A\subset {\Bbb R}^n$
 و $َA$ مستقل خطی باشد آنگاه :
 $span(A)={\Bbb R}^n$.
 \\
 د) اگر 
 $span(A)={\Bbb R}^n$
 آنگاه $A$ یک مجموعه مستقل خطی از
 ${\Bbb R}^n$
 است.
 \\
 ه) اگر 
 $w$
 یک ترکیب خطی از بردار های 
 $v_1,v_2,\ldots,v_n$
 باشد،آنگاه 
 $\{w,v_1,v_2,\ldots,v_n\}$
 مستقل خطی است.
 \\
 ی)اگر 
 $S \subseteq {\Bbb R}^n$
 مسنقل خطی باشد و 
 $v \in ({{\Bbb R}^n}-span(S))$
 آنگاه 
 $S\cup \{v\}$
 نیز مستقل خطی است.
 
 \مسئله{} مانریس سودوکو یک ماتریس 
 $9\times 9$
 که اعداد 
 $1,2,\ldots,9$
در هر سطر در هر ستون و هر بلوک 
$9\times 9$
آن ظاهر شده اند ، اگر $S$یک ماتریس سودوکو باشد به سوالات زیر پاسخ دهید:
\\
الف) حاصلضرب 
$S {\left(\begin{array}{ccc}
	1\\
	1\\
	\vdots\\
	1
\end{array}\right)_{9\times 1}}$
 را بیابید.
 \\
 ب)  کدام یک از اعمال سطری روی 
 $S$
 ،مجددا یک ماتریس سودوکو به ما می دهد؟
 
 \مسئله {} فرض کنید 
 $ََA$
یک ماتریس 
$m\times n$
است که:
\\
الف) برای هر $b$  در 
${\Bbb R}^n$
 معادله 
 $Ax=b$
 حداکثر یک جواب دارد،ثابت کنید ستون های ماتریس 
 $A$
 باید مستقل خطی باشند.
 \\
 ب) که $n$ تا از ستون های آن محوری هستند،ثابت کنید برای هر $b$  در 
 ${\Bbb R}^n$
 معادله ${\Bbb R}^n$ حداکثر یک جواب دارد.
 \newpage
 \مسئله {} مشخص کنید هریک از تبدیلات زیر خطی هستند یا نه،در صورتی که خطی باشند ماتریس استاندارد آن ها را نیز بیابید.\\
 الف)
 {\setlength\arraycolsep{0.1em}
 \begin{eqnarray*}
 	f:{\Bbb R}^2&\longrightarrow& {\Bbb R}^3\\
 	 (x,y)&\longmapsto&(x^2,2y)
 \end{eqnarray*}}
\\
ب) 
{\setlength\arraycolsep{0.1em}
	\begin{eqnarray*}
		f:{\Bbb R}^2\ \ &\longrightarrow&\ \ {\Bbb R}^2\\
		(x,y)&\longmapsto&(2x+y,-y)
\end{eqnarray*}}\\
ج) اگر تبدیل خطی زیر به شکل:
{\setlength\arraycolsep{0.1em}
	\begin{eqnarray*}
		f:{\Bbb R}^2\ \ &\longrightarrow&\ \ {\Bbb R}^2\\
		(v_1,v_2)&\longmapsto&(\frac{v_1+v_2}{2},\frac{v_1+v_2}{2})
\end{eqnarray*}}
آنگاه در مورد 
$f(f(v_1,v_2))$
چه می توان گفت؟(مشخص کنید خطی هست یا نه و درصورت خطی بودن ماتریس استاندارد آن را بیابید.)
\end{document}
